\documentclass[11pt,oneside]{article}
\frenchspacing

\usepackage[utf8]{inputenc}
\usepackage[T1]{fontenc}
\usepackage{amsmath}
\usepackage{amssymb}
\usepackage{amsthm}
\usepackage{enumitem}
\usepackage{geometry}
\usepackage{hyperref}

\geometry{margin=1in}
\hypersetup{
  colorlinks=true,
  linkcolor=blue,
  urlcolor=blue,
}

\newtheorem{definition}{Definition}
\newtheorem{proposition}{Proposition}

\newcommand{\R}{\mathbb{R}}
\newcommand{\barp}{\bar p}
\newcommand{\estate}{e}

\title{Eisenberg--Noe Clearing as a Smart Contract Primitive}
\author{Research note}
\date{June 2026}

\begin{document}
\maketitle

\begin{abstract}
The research idea is to implement a self-clearing credit network on-chain.
Borrowers and lenders, or pools thereof in lending architectures such as
Morpho, Aave, and related venues, do not resolve distress through isolated
liquidations or first-come-first-served withdrawals. They opt into a common
liability graph.
When the network enters settlement, a smart contract applies an
Eisenberg--Noe style global bankruptcy rule: each debtor pays what it can, each
creditor receives pro rata within its claim class, and circular obligations are
cleared by a fixed point. The purpose of the contract is not to make default
disappear. It is to replace local liquidation races with a verifiable global
clearing vector.
\end{abstract}

\section{Motivation}

Most DeFi credit protocols are pairwise. A position becomes unsafe, a liquidator
buys collateral, and the local loan is closed. This is clean for a single
overcollateralized loan, but it is less natural for a network of vaults,
borrowers, lenders, rehypothecated claims, recursive lending loops, or tokenized
credit lines.

In such a network, a participant's solvency may depend on payments from other
participants who themselves depend on payments from the first participant or
from a shared set of intermediaries. Local liquidation can then destroy value
or amplify a stress that a global settlement would have absorbed.

The Eisenberg--Noe model gives a compact rule for this setting
\cite{eisenberg2001systemic}. It treats the financial system as a directed
graph of nominal obligations and computes the clearing payments that are
mutually consistent with limited liability and pro rata repayment. The
smart-contract question is:

\begin{quote}
Can we build an on-chain clearinghouse where the liability graph is native
state, the clearing vector is proposed by solvers, and the contract verifies
and executes the fixed point?
\end{quote}

\section{The Clearing Model}

Let $N=\{1,\ldots,n\}$ be a set of participants. Participant $i$ has external
assets $x_i \geq 0$ and nominal liabilities $\barp_i \geq 0$. The matrix
$\Pi \in \R_+^{n \times n}$ records relative liabilities: $\Pi_{ij}$ is the
fraction of $i$'s total liabilities owed to $j$. If $\barp_i>0$ then
\[
  \sum_j \Pi_{ij}=1.
\]
If $\barp_i=0$, set $\Pi_{ij}=0$ for all $j$.

A clearing payment vector $p \in \R_+^n$ satisfies
\[
  0 \leq p_i \leq \barp_i
\]
and
\[
  p_i =
  \min\left\{
    \barp_i,\,
    x_i + \sum_j \Pi_{ji}p_j
  \right\}.
\]
The term
\[
  \estate_i(p) = x_i + \sum_j \Pi_{ji}p_j
\]
is the estate available to participant $i$ after incoming payments. If
$\estate_i(p) \geq \barp_i$, participant $i$ pays all liabilities. If
$\estate_i(p)<\barp_i$, participant $i$ defaults and pays exactly its estate
pro rata to creditors.

Equivalently, define the monotone map
\[
  \Phi_i(p)=
  \min\left\{
    \barp_i,\,
    x_i + \sum_j \Pi_{ji}p_j
  \right\}.
\]
A clearing vector is a fixed point $p=\Phi(p)$. Under standard regularity
conditions the relevant clearing vector is unique. Recent work sharpens the
equilibrium uniqueness question for the original model
\cite{stachurski2022properties}. Even when uniqueness needs care, the greatest
clearing vector has a clear economic meaning: it is the most optimistic payment
vector compatible with limited liability and the liability graph.

\section{Contract Interpretation}

The on-chain primitive is a settlement contract for a single numeraire asset at
first. Participants deposit assets, issue obligations to each other, and agree
that a settlement epoch will be resolved by the clearing rule.

\subsection{State}

The minimal state is:

\begin{itemize}[leftmargin=1.5em]
  \item a participant registry;
  \item a set of directed liability edges $(i,j,L_{ij})$;
  \item deposited external assets $x_i$ for each participant;
  \item a settlement status: live, frozen, proposed, challenged, or executed;
  \item the proposed clearing vector $p$ for the current epoch.
\end{itemize}

From the raw liability matrix $L$, the contract derives
\[
  \barp_i=\sum_j L_{ij}
  \quad\text{and}\quad
  \Pi_{ij} =
  \begin{cases}
    L_{ij}/\barp_i, & \barp_i>0,\\
    0, & \barp_i=0.
  \end{cases}
\]

\subsection{Lifecycle}

A simple lifecycle is:

\begin{enumerate}[leftmargin=1.5em]
  \item Live mode. Participants create, transfer, repay, or cancel obligations
  subject to protocol rules.
  \item Freeze. An epoch ends or a distress trigger freezes the graph.
  \item Proposal. A solver submits a clearing vector $p$ and posts a bond.
  \item Verification. The contract checks that $p$ satisfies the clearing
  equations within exact integer arithmetic or a stated tolerance.
  \item Challenge. Other participants can challenge an invalid vector during a
  short window.
  \item Execution. The contract distributes assets according to $p$ and marks
  unpaid liabilities as defaulted, rolled, tokenized, or extinguished according
  to the product design.
\end{enumerate}

\section{Why Verification Is Easier Than Solving}

The contract does not need to run an iterative fixed-point algorithm for large
graphs. A solver can compute $p$ off-chain. The contract only needs to verify
that the proposed vector is a fixed point.

For each participant $i$, compute
\[
  \estate_i = x_i + \sum_j L_{ji}\frac{p_j}{\barp_j},
\]
where the fraction is defined as zero if $\barp_j=0$. Then check:
\[
  0 \leq p_i \leq \barp_i.
\]
If $p_i < \barp_i$, require
\[
  p_i = \estate_i.
\]
If $p_i=\barp_i$, require
\[
  \estate_i \geq \barp_i.
\]
These local checks prove the global fixed-point condition. They cost
$O(n+|E|)$ for a graph with $n$ participants and $|E|$ liability edges.

\subsection{Certifying the greatest fixed point}

The canonical Eisenberg--Noe clearing vector is the greatest fixed point of
\[
  \Phi(p)=\barp \wedge (x+\Pi^\top p).
\]
Economically, it is the maximal payment vector compatible with limited
liability, priority, and the liability graph. If the clearing vector is unique,
then any verified fixed point is already canonical. If uniqueness is not built
into the admissible graph, verifying
\[
  p=\Phi(p)
\]
only proves that $p$ is a fixed point. It does not by itself prove that $p$ is
the greatest one.

The direct certificate is the top-iteration trace:
\[
  p^0=\barp,\qquad p^{t+1}=\Phi(p^t).
\]
Since $\Phi$ is monotone and $\Phi(\barp)\leq \barp$, the sequence decreases.
In exact integer arithmetic with a deterministic rounding rule, the trace
terminates when
\[
  p^{T+1}=p^T.
\]
The terminal vector is the greatest fixed point. A solver can therefore submit
the whole trace
\[
  p^0,p^1,\ldots,p^T
\]
and the contract verifies every transition. This is stronger than a local
fixed-point check, but it costs $O(T|E|)$ rather than $O(|E|)$.

There are several lighter variants:

\begin{itemize}[leftmargin=1.5em]
  \item enforce uniqueness conditions on admissible graphs, so a local
  fixed-point check is sufficient;
  \item use an optimistic challenge rule where any participant can disprove a
  proposal by submitting a strictly higher valid fixed point;
  \item require a proof that the submitted vector is the result of the
  top-iteration procedure, possibly with a succinct proof system;
  \item restrict the first implementation to graph classes where multiplicity is
  not the core risk.
\end{itemize}

This maximality issue may become less central in a soft-clearing design. If the
system moves through warning, compression, and partial deleveraging before hard
settlement, the protocol may care more about conservative executable recovery
values and monotone risk triggers than about proving a terminal bankruptcy
allocation in one step. Still, for a hard clearing epoch, the contract should
either verify the greatest fixed point or make explicit why any fixed point is
canonical.

\subsection{Integer Form}

In Solidity, the contract should avoid rational numbers. It can store raw
liabilities $L_{ij}$ and payments $p_i$ in integer token units. To check
participant $i$, multiply through by common denominators or use conservative
rounding rules.

For each incoming edge $j \to i$, the payment sent from $j$ to $i$ is
\[
  q_{ji} = \left\lfloor \frac{L_{ji}p_j}{\barp_j} \right\rfloor.
\]
The rounding remainder can be assigned to a deterministic residue account, sent
to creditors by edge order, or retained as dust. The important design point is
that the rounding rule must be part of the clearing equation, not an
implementation afterthought.

\section{Smart Contract Sketch}

The verifier has this shape:

\begin{verbatim}
for each participant i:
    require(p[i] <= totalLiabilities[i])
    estate = externalAssets[i]

for each liability edge j -> i with amount L[j][i]:
    if totalLiabilities[j] > 0:
        estate[i] += floor(L[j][i] * p[j] / totalLiabilities[j])

for each participant i:
    if p[i] < totalLiabilities[i]:
        require(p[i] == estate[i])
    else:
        require(estate[i] >= totalLiabilities[i])
\end{verbatim}

Execution then pushes edge payments:

\begin{verbatim}
for each liability edge i -> j with amount L[i][j]:
    if totalLiabilities[i] > 0:
        q = floor(L[i][j] * p[i] / totalLiabilities[i])
        transfer(asset, j, q)
\end{verbatim}

This is the core self-solving mechanism. Solvers compete to submit the correct
clearing vector. The contract enforces the bankruptcy rule. Participants no
longer need to race each other through local closeouts once the network is
frozen.

\section{What This Buys}

\subsection{Netting circular claims}

Circular obligations are not exceptional in the model. If $A$ owes $B$, $B$
owes $C$, and $C$ owes $A$, the fixed point clears the cycle jointly. The result
can preserve value that pairwise liquidation would waste.

\subsection{Contagion as state}

The default set is endogenous:
\[
  D(p)=\{i: p_i<\barp_i\}.
\]
Losses, contagion paths, and recovery rates become explicit outputs of the
settlement contract. This is useful for monitoring, insurance, tranche design,
and risk capital.

\subsection{Bankruptcy before liquidation}

The contract separates the insolvency rule from the asset liquidation rule.
In a collateralized DeFi market, collateral can be sold first and residual
unsecured claims can then enter the clearing graph. Or collateral can remain in
the estate and be distributed by the global clearing vector. These are different
products, but the same fixed-point logic can sit beneath both.

\section{DeFi Extensions}

The plain Eisenberg--Noe model assumes a single numeraire, fixed nominal
liabilities, and pro rata claims of equal seniority. DeFi needs extensions.
Close relatives in the financial-network literature add default costs, rescue
rules, liquidity effects, and network-stability analysis
\cite{rogers2013failure,cifuentes2005liquidity,acemoglu2015systemic,glasserman2016contagion}.

\subsection{Collateral}

Collateral can be handled by a waterfall:

\begin{enumerate}[leftmargin=1.5em]
  \item seize or mark collateral;
  \item repay secured creditors up to the collateral value;
  \item place any residual deficiency into the unsecured clearing graph;
  \item clear unsecured liabilities by the Eisenberg--Noe rule.
\end{enumerate}

An alternative is to treat collateral as part of $x_i$ and allow the global
clearing to distribute the estate. That is simpler but may conflict with
secured-creditor expectations.

\subsection{Seniority}

Multiple seniority classes can be implemented by sequential clearing. Senior
claims are cleared first. Remaining estate then enters the next class. This
preserves priority without abandoning the network view.

\subsection{Multiple assets}

A multi-asset version requires prices or exchange mechanisms. The cleanest MVP
uses a single ERC20 numeraire. A later version can introduce:

\begin{itemize}[leftmargin=1.5em]
  \item oracle prices fixed at the freeze block;
  \item on-chain auctions for non-numeraire collateral;
  \item separate clearing vectors by asset;
  \item endogenous price impact if liquidation size affects recovery.
\end{itemize}

\subsection{Dynamic obligations}

The graph must freeze before settlement. Otherwise participants can create
claims after observing the proposed clearing vector. A real implementation needs
epoch boundaries, rate limits, permissioning, or a mempool-resistant freeze
rule.

\section{Objections as Design Constraints}

The natural objections narrow to two design questions:

\begin{itemize}[leftmargin=1.5em]
  \item which claims are eligible for the clearing graph;
  \item how secured and collateralized claims enter the clearing rule.
\end{itemize}

Some apparent objections are implementation choices rather than blockers.
Solving the fixed point on-chain is unnecessary: the solver can compute the
candidate vector off-chain, while the contract verifies the fixed-point
conditions in $O(n+|E|)$. The remaining subtlety is maximality: if uniqueness
conditions are not enforced by the admissible graph, the contract should require
the greatest clearing vector rather than merely any fixed point, using one of
the certificate or challenge rules above.

Oracle manipulation is also not special to this construction. Any DeFi system
that admits marked collateral inherits oracle risk. The sharper version is that
the oracle becomes a bankruptcy allocation input. The first prototype should
therefore use a single deposited numeraire and no oracle. Marked collateral can
enter only after the fixed-point verifier and settlement accounting are clean.

\subsection{Why plain pro rata bankruptcy is not enough}

Plain Eisenberg--Noe is pari passu unsecured bankruptcy. Each creditor of a
defaulting node shares the estate pro rata within a claim class. DeFi lending is
usually not shaped like that. It has secured claims, liquidation thresholds,
collateral-specific liens, seniority, liquidation penalties, and sometimes
price impact.

Consider a debtor $A$ with:

\[
\begin{array}{rcl}
  A & \text{owes} & B: 100, \quad \text{secured by collateral worth } 80,\\
  A & \text{owes} & C: 100, \quad \text{unsecured},\\
  A & \text{has} & 20 \quad \text{free assets}.
\end{array}
\]

Plain pro rata clearing would see estate $100$ and liabilities $200$, so $B$
receives $50$ and $C$ receives $50$. That ignores the secured nature of $B$'s
claim.

A secured-credit waterfall gives a different result. First, $B$ receives the
collateral value $80$. This leaves a residual unsecured deficiency of $20$ for
$B$. The unsecured graph then has claims $B:20$ and $C:100$ against free estate
$20$, so $B$ receives $20\cdot 20/120=3.33$ and $C$ receives
$20\cdot 100/120=16.67$. The final allocation is:

\[
  B:83.33,\qquad C:16.67.
\]

This suggests the clean extension:

\begin{enumerate}[leftmargin=1.5em]
  \item resolve secured closeout or collateral waterfalls first;
  \item convert collateral shortfalls into unsecured deficiency claims;
  \item add any collateral surplus back to the debtor estate;
  \item run Eisenberg--Noe on the residual unsecured graph.
\end{enumerate}

The fixed-point logic is then preserved, but it applies to the residual claims
that should actually be pari passu.

\subsection{Strategic claim creation near clearing}

The strongest strategic issue is not hidden bilateral obligations. Hidden claims
should simply be ineligible for the current clearing. The dangerous case is a
visible but late obligation or claim transfer inserted once participants know
that clearing is near.

Examples include:

\begin{itemize}[leftmargin=1.5em]
  \item a distressed borrower creates a friendly liability to an affiliated
  address;
  \item a creditor transfers claims to a party with different treatment;
  \item claims are routed through an entity expected to default in order to
  change the loss allocation;
  \item a participant front-runs the freeze trigger with a new edge.
\end{itemize}

The defense is an eligibility and seasoning rule:

\begin{quote}
Only obligations created before snapshot block $B-k$ enter the current
clearing. Late obligations roll into the next epoch. Claim transfers preserve
the original claim age and seniority. New obligations must arise from
protocol-verified value transfer, not arbitrary IOU creation.
\end{quote}

This makes the clearing graph a state machine rather than a free-form list of
asserted debts.

\subsection{Soft deleveraging before hard clearing}

The network does not need to jump directly from normal lending to bankruptcy
settlement. Clearing can be soft and pre-emptive:

\begin{center}
\begin{tabular}{ll}
normal mode & ordinary borrowing, lending, repayment, and transfers,\\
warning mode & risk limits tighten and new concentration is discouraged,\\
compression mode & new debt is disabled and repayments remain open,\\
freeze mode & the eligible graph is snapshotted,\\
clearing mode & the clearing vector is verified and executed.
\end{tabular}
\end{center}

Triggers can be scheduled or risk-based: projected default under a haircut,
network leverage above a bound, collateral liquidity below a bound, oracle
deviation, failed margin top-up, participant petition, or governance call.

\subsection{A candidate clearing rule}

A concrete first rule is:

\begin{enumerate}[leftmargin=1.5em]
  \item Snapshot the liability graph and admitted collateral at block $B$.
  \item Exclude non-seasoned claims from the current epoch.
  \item Close secured claims against their specific collateral first.
  \item Convert collateral shortfalls into unsecured deficiency claims.
  \item Add collateral surplus back to the debtor estate.
  \item Run Eisenberg--Noe on the residual unsecured graph.
  \item Execute pro rata payments and, if useful, tokenize unpaid residual
  claims.
\end{enumerate}

The rule keeps collateral priority explicit while still using global clearing
for the part of the system where global clearing is the correct object:
unsecured residual claims and circular payment dependencies.

\section{Research Questions}

\begin{itemize}[leftmargin=1.5em]
  \item What is the right trigger: scheduled epoch, solvency breach, governance
  call, oracle event, or participant petition?
  \item How should solver bonds and challenges be designed so invalid vectors
  are cheap to reject and valid vectors are cheap to finalize?
  \item Can the verifier be compressed with Merkle commitments, fraud proofs, or
  zero-knowledge proofs for large graphs?
  \item How do we handle collateral priority without reducing the system back to
  isolated liquidations?
  \item Which DeFi liabilities are eligible: vault shares, debt tokens, margin
  accounts, credit lines, receivables, or RWA notes?
  \item What strategic behavior appears when agents know that obligations will
  be globally cleared?
  \item Can the clearing outputs feed insurance premia, capital buffers, or
  curator risk limits?
\end{itemize}

\section{Minimal Viable Experiment}

The first useful implementation should be deliberately narrow:

\begin{itemize}[leftmargin=1.5em]
  \item single ERC20 numeraire;
  \item no oracle;
  \item all external assets are deposited in the clearing contract;
  \item fixed participant set during each epoch;
  \item raw liability edges registered on-chain;
  \item off-chain solver;
  \item on-chain fixed-point verification;
  \item optimistic challenge window;
  \item deterministic integer rounding.
\end{itemize}

This version is enough to test the central claim: a smart contract can enforce a
global bankruptcy rule over a liability graph without solving the fixed point
on-chain.

\section{Connection to Recursive Lending}

The model is especially relevant to recursive vault and lending structures.
When a vault lends against its own shares, or when vault shares become collateral
inside another lending venue, the dependency graph can contain cycles. Pairwise
liquidation sees only the local edge. A global clearing contract sees the cycle.

This suggests a product direction:

\begin{quote}
Build a credit network where recursive DeFi positions opt into a common
settlement layer. Normal times use ordinary lending mechanics. Stress times use
global clearing.
\end{quote}

The resulting system would sit between two familiar designs. It is not a simple
AMM or isolated lending market. It is also not a discretionary bankruptcy court.
It is a rule-based clearing court implemented as a contract.

\section{Near-Term To-Dos}

\begin{itemize}[leftmargin=1.5em]
  \item Write the single-numeraire contract invariant precisely.
  \item Work through a three-node cycle example by hand.
  \item Specify the integer rounding rule and dust ownership.
  \item Decide whether the first prototype should execute payments directly or
  mint recovery tokens.
  \item Compare optimistic verification, full on-chain verification, and zk
  verification.
  \item Connect the graph input to YTD or DIG-style dependency extraction.
\end{itemize}

\section{Varia}

For unwinding thresholds, the relevant asset value is not necessarily
mark-to-market value. A clearing network should consider computing
liquidative values: what each position can actually realize when it must be
unwound through the available exit routes.

This matters because a position can look solvent at quoted prices while being
insolvent at executable exit depth. The clearing trigger should therefore be
able to use haircut or liquidation-value estates,
\[
  x_i^{\mathrm{liq}} \leq x_i^{\mathrm{mtm}},
\]
where $x_i^{\mathrm{liq}}$ reflects pool depth, redemption queues, auction
discounts, oracle delay, venue concentration, and expected slippage under the
size of the unwind. In that version, the warning and compression modes are not
triggered by accounting insolvency alone. They are triggered when the network
approaches the point where executable liquidation value no longer covers the
eligible claims.

\bibliographystyle{plain}
\bibliography{Eisenberg-Noe}

\end{document}
